\chapter{Research design and evaluation}
\label{ch:data}

The evaluation follows three questions: whether VolPFN forecasts usefully, what its training choices contribute, and when correcting the conversion from scale to tails helps. Table~\ref{tab:study_map} identifies the model and sample used for each study. We call one experimental configuration an \emph{arm}, and a saved set of trained weights a \emph{checkpoint}. The full configuration includes extra supervision from simulated volatility labels; \emph{noaux} removes that auxiliary loss. Direct predicts expected squared return, while historical learns a return distribution from market windows. The following sections specify these procedures and the comparisons they support.

\section{Questions and completed comparisons}

For each study, we specify the forecast object, the model issuing it and the criterion for success. A \emph{seed} controls pseudorandom training choices. Five runs allow us to check sensitivity to initial weights and simulated tasks, although they share the market observations used for evaluation \citep{bouthillier2021variance}. They therefore provide a training-sensitivity check rather than five independent market samples.

\begin{longtable}{@{}>{\raggedright\arraybackslash}p{2.7cm}>{\raggedright\arraybackslash}p{5.5cm}>{\raggedright\arraybackslash}p{6.4cm}@{}}
\caption[Map of the completed studies]{Completed studies and their separate roles. A market--date row is a scored observation, not an independent replication. A union date is observed in at least one included market.}\label{tab:study_map}\\
\toprule Study & Population and endpoint & Question and model identity \\
\midrule\endfirsthead
\toprule Study & Population and endpoint & Question and model identity \\
\midrule\endhead
Corrected foundation benchmark & 28 indices; 8,774 rows; 327 union dates; 1 April 2015--30 June 2016. Development-scaled $r^2$ QLIKE. & Five corrected broad models against three fixed pretrained adapters; historical confirmation. \\
Matched prior & 20 indices; 22,541 rows; 1,176 union dates; 1 March 2022--31 August 2026. Raw M2, common-$t_8$ log score. & Five matched broad--Mat\'ern pairs; prior-package contribution. \\
Components and financial neural models & 28 markets; 50,073 return rows per seed; 1 January 2015--25 February 2022. Raw $r^2$ QLIKE; secondary RK mask 50,054. & Four configurations at five seeds each, plus three dedicated daily adaptations; distinct three-/four-slot families. \\
Harmonised daily & 28 markets; 22,026 keys; 827 union dates; 1 January 2019--25 February 2022. Raw M2/$r^2$ and variance/RK QLIKE. & Saved forecasts aligned without retraining or new calibration; full 30 September reference; 28-slot family. \\
Intraday extension & Same 22,026 keys and dates; two proxy-specific procedures. & Frozen daily VolPFN+TTM reference, stronger realised-risk adaptations and matched GJR blends; 38-slot family. \\
Cross-market helper & Same 22,026 keys and dates; five helper seeds. & VolPFN+helper reference, calibration and GJR controls; separate twelve-slot follow-up. \\
Historical follow-up & Same 28 indices and 50,073 return rows per seed; 2015--25 February 2022. Raw $r^2$ QLIKE. & Five regularised, development-selected historical models versus saved full/noaux; separate two-slot family. \\
TTM calibration controls & Same 22,026 keys, 28 markets and 827 union dates; 2019--February 2022. Separate $r^2$/RK QLIKE. & Saved full+TTM blend versus affine and scale TTM-only controls; separate four-slot family. \\
External ETF population & 20 fixed ETFs; 22,600 rows per model/seed; 1 March 2022--31 August 2026. Raw $r^2$ QLIKE. & Frozen full/noaux and index-selected historical weights versus rolling GJR; separate three-slot family. \\
Historical explanation and tails & Development 2011--2014; later 2015--February 2022; original masks and scores. & Retained checkpoint, independently fitted compact rule and archived tail procedures; repaired sensitivity stays separate. \\
\bottomrule\end{longtable}

The study map also identifies boundaries between results. Reusing seed numbers does not make different checkpoints interchangeable, and studies with different periods or outputs cannot be pooled without qualification. In particular, subtracting a prior gain from a benchmark margin measured elsewhere would not explain that margin.

\section{Target, timing and available history}

We retain the archive's target, $100\log(C_t/O_t)$, where $O_t$ is opening price and $C_t$ is closing price. This gives the session log return in percent units, with additive log changes, and keeps the programme's studies comparable. Because the interval excludes overnight exposure, the evaluation does not measure continuously held portfolio performance or trading profits.

A \emph{forecast origin} is the point at which the prediction is issued. Figure~\ref{fig:forecast_origin} shows why all inputs must precede it. The same restriction applies to scaling: using future observations to normalise the history introduces leakage even with causal attention. Both preprocessing and prediction must therefore preserve chronological information availability \citep{tashman2000out,kapoor2023leakage}.

\begin{figure}[htbp]
\centering
\begin{tikzpicture}[x=1cm,y=1cm,>={Stealth},every node/.style={font=\small,align=center}]
\fill[blue!6] (0,-0.43) rectangle (7.3,0.43);
\fill[orange!10] (8.3,-0.43) rectangle (12.9,0.43);
\draw[->,thick] (0,0)--(13.4,0);
\foreach \x in {1,2.4,3.8,5.2,6.6}{\draw (\x,-0.1)--(\x,0.1);}
\draw[thick,dashed] (8,-0.85)--(8,1.2);
\node[text width=6.6cm] at (3.6,1.05) {\textbf{Observed history}\\Past completed returns only};
\node[text width=4.6cm] at (10.8,1.05) {\textbf{Target session}\\Opening to closing return};
\node[text width=6.5cm] at (3.6,-1.1) {Compute features and scale\\using information already available};
\node[text width=4.5cm] at (10.8,-1.1) {Observe the return after the session,\\then evaluate the saved forecast};
\node[font=\small\bfseries,text width=2.7cm] at (8,1.9) {Forecast issued};
\draw[->] (8,1.57)--(8,1.22);
\end{tikzpicture}
\caption[Forecast timing and information availability]{A schematic forecast origin. The target session is not part of the input history. The drawing explains information order; it is not a plot of observed prices or an exchange-clock model. Historical benchmark forecasts use 1,024 preceding observed returns at each origin.}
\label{fig:forecast_origin}
\end{figure}

The historical benchmarks receive 1,024 preceding observed returns, with non-trading days omitted. The archive fixes this history length; its optimality has not been demonstrated. The prior study uses a different interface: windows of at most 1,024 positions, a stride of 512 and at least 100 earlier observations. Where windows overlap, a fixed ownership rule assigns each prediction once. Both prior arms use that rule, preserving their match without equating their processing interface to the historical benchmark's.

\section{Scores chosen for the forecast object}

The score must correspond to the forecast object. Risk scoring evaluates the second moment, while density scoring evaluates the complete distribution. Each study has a designated primary score, which we retain when reporting its result. The timing of the analysis decisions matters separately: the exposure record below distinguishes plans made before scoring from retrospective specifications.

\paragraph{Risk scale.}
We assess a positive risk forecast $h$ through QLIKE, $r^2/h+\log h$, with lower loss preferred. Squared return noisily observes the second moment and equals variance in expectation under zero mean. Under the required assumptions, this proxy noise need not distort expected rankings \citep{patton2011volatility}. We report competitor loss minus VolPFN loss as gain, so a positive value favours VolPFN. An absolute gain of 0.01 does not automatically mean a one-percent improvement.

\paragraph{Predictive distributions.}
For a complete distribution, we score the log density assigned to the observed return. This proper distributional score prefers larger values \citep{gneiting2007strictly}. Natural logarithms give units of \emph{nats}: a gain of 0.005 corresponds to a geometric-mean density ratio $\exp(0.005)$, about 1.005. The roughly half-percent increase describes density at observed outcomes, rather than investment return.

For the prior comparison, we convert both second-moment forecasts to the same zero-centred Student distribution with eight degrees of freedom, termed the \emph{common-$t_8$ law}. Holding this scoring shape fixed prevents a tail-law difference from being interpreted as a scale gain. We use eight as a shared convention; it was not estimated as an optimum. The network's output head still determines the mixture second moment, labelled M2, used in the conversion. Supporting \emph{native-density} scores instead evaluate the distribution each model actually issues.

\section{Data qualification before scoring}

The original daily data comprise 28 Oxford--Man indices from Barde's replication material \citep{heber2009oxford,barde2024march}. Both the original and added daily procedures receive returns alone. A realised kernel (RK) estimates intraday variability. We use it as an alternative risk proxy or, where specified, a lagged input. It is not observed conditional variance, and its conditional expectation has not been equated to next-session return risk. Consequently, the interpretation of squared-return QLIKE does not automatically extend to RK rankings.

For later histories, we map the same 28 index identities to named Yahoo indices before acquisition and retain raw responses, metadata and prepared open--high--low--close files. A frozen \emph{scoring mask} determines eligible observations before forecast errors are inspected. The recording checks address data-quality concerns discussed by \citet{ince2006data}, while their thresholds remain our protocol choices rather than requirements of that paper.

A market is eligible with 800 valid evaluation observations, 100 preceding observations, no more than 1\% invalid primary rows, and unique increasing dates. Comparable opening prices may repeat the preceding close in at most 20\% of cases within tolerance. Frequent repetition raises a question about separating session and overnight movement, but rounding and valuation conventions can also cause it. The check therefore identifies a recording concern without proving that the prices are wrong.

The fixed gate admits twenty markets and 22,541 observations. These cutoffs are neither demonstrated optima nor a power calculation, and passing them does not independently verify the prices. We retain the excluded markets without replacement or target changes. Appendix~\ref{app:closure_protocol} records all eight exclusion reasons and the relevant tolerances.

The external follow-up freezes twenty ETF identities before acquiring Yahoo histories. All pass the recorded rules and contribute 1,130 evaluation sessions each in March 2022--August 2026. The same session-return target uses unadjusted opening and closing prices, vendor-local dates and retained zero returns. This extends the asset population, but supplies neither an independent calendar holdout nor another broad--Mat\'ern comparison. Appendix~\ref{sec:followup-methods} gives the roster and source policy.

\section{The component and financial neural programme}
\label{sec:new-programme-design}

We use four configurations to examine particular training choices. \emph{Full} combines the broad prior and return mixture with auxiliary log-variance supervision weighted 0.1. \emph{Noaux}, meaning no auxiliary loss, removes that objective while retaining the unused head and matched initialisation. \emph{Direct} shares the initial trunk and synthetic batches but predicts a positive second moment through squared-return QLIKE. \emph{Historical} applies the noaux density objective to market windows. The direct comparison changes output and loss together; the historical comparison changes the source and diversity of tasks.

Each configuration uses seeds 8101--8105, with 20,000 updates of 8,192 positions per run, giving 163,840,000 positions per model. All twenty runs completed. We use a \emph{fixed-final} checkpoint at the final scheduled update, here 20,000, without validation-based weight selection. Historical windows end on 31 December 2010. Although supplied compute matches, unique information does not: historical observations repeat, whereas simulation supplies new tasks. The equal budget therefore controls the number of training positions without making the two sources equally diverse.

The dedicated financial models fit each market through 2010, select settings on 2011--2014 and evaluate 1 January 2015--25 February 2022. ARCH-LSTM, GARCH-LSTM and GINN receive contexts of one, 126 and ninety observations respectively. Means and centred root-mean-square (RMS) scales are fitted only on training observations, denoted \emph{TRAIN}, and then frozen. Their grids allow fifty epochs, patience eight and batches of 128, with development squared-return QLIKE selecting settings (Appendix~\ref{sec:new-research-methods}). This gives them selection opportunities absent from fixed-final simulation training.

Evaluation retains five seeds, all 28 markets and 50,073 returns per seed, including zeros. Requiring finite positive RK gives a separate 50,054-observation target mask. A failure cannot remove difficult observations from the population. Unavailable RECH retains its comparison slot at $p=1$, while DCC/GN-GARCH was not run and cannot be described as defeated.

\paragraph{A stronger historical comparison.}
The retrospective historical follow-up strengthens the comparator by adding regularisation and allowing earlier development observations to select its weights. It retains the base mixture and noaux density objective, with residual dropout 0.1. We train the model on pre-2011 index windows using AdamW, a learning rate of $10^{-4}$, weight decay 0.1 and 8,192 positions per update.

Selection uses equal-market squared-return QLIKE on 1,792 fixed development origins from 2011--2014. At each 500-update check, we retain the first strict minimum; six checks without improvement stop training, subject to a 20,000-update cap. Every seed selects update 500 and stops at 3,500 before scoring the original 50,073 observations per seed. These selected historical weights therefore receive fewer updates than the fixed-final synthetic weights. Because source, dropout and selection also differ, the follow-up compares training procedures rather than isolating source alone.

The ETF study reuses the improved historical checkpoints and original full/noaux weights without ETF training or checkpoint selection. Its rolling GJR-GARCH-$t$ comparator fits a constant mean on 1,000 preceding returns and refits every twenty market sessions. Variance plus squared fitted mean supplies its squared-return forecast. Four of 1,140 refit attempts fail; each carries forward the preceding accepted parameters and continues daily filtering, retaining every evaluation observation. Appendix~\ref{sec:followup-methods} specifies this policy and the causal calibration controls.

\section{Matching learning opportunity}

Matching the prior configurations requires architecture, transformations, losses, optimiser, precision and position allocation to agree \citep{dodge2019show}. Both receive the 20,000 updates specified above. The supplied-position total includes targets masked out of the loss, and equal allocations alone do not establish convergence.

Contexts have lengths 512 and 1,024 with probabilities 0.4 and 0.6, giving more weight to the longer evaluation context. We match positions and updates rather than elapsed time, so a cheaper simulator gains no additional training. Return innovations and observation modifications also match, leaving the combined latent-volatility design as the intervention.

Within seeds 8101--8105, each pair starts from identical initial states. Seed and update index determine batch generation independently of scheduling. We retain update 20,000 rather than selecting weights from market outcomes or a validation mixture favouring one prior; checks every 500 updates monitor training health. The primary comparison requires all ten final models. Eight training runs are new; two final Mat\'ern checkpoints are audited reuses, as recorded in Appendix~\ref{app:closure_protocol}.

\section{Preprocessing and comparison interfaces}

\subsection{Causal scale and features}

A running RMS based on preceding nonzero observations brings market magnitudes onto comparable scales while preserving squared-return units in the forecast. Before any nonzero observation, the scale is one. The next target is always transformed with the scale available before it occurs, keeping the transformation causal.

Exponential moments skip zero returns under the simulator's stale-price convention. This is a modelling choice because genuine market zeros also occur. Training masks stale targets and uses a continuous mixture with no mass concentrated exactly at zero, so matching market masks does not establish a zero-inclusive predictive law. We have not isolated alternative RMS choices or the benefit of window resets. Resetting aligns contexts while limiting memory.

The risk summaries use half-lives of five, twenty and 100 nonzero observations, motivated by heterogeneous horizons and exponential weighting \citep{corsi2009simple,riskmetrics1996}. We divide their raw-unit moments by the current squared RMS (Appendix~\ref{app:forecast-details}). They summarise daily returns, which distinguishes them from HAR's intraday inputs.

\subsection{Adapting general forecasters}

An \emph{adapter} converts a model's inputs and outputs to the scored forecast. Chronos-2 and Moirai-2 take squared-return histories and expose medians, while TabPFN-TS requests a mean with eight causal covariates and a running index. Information availability matches, but representations and forecast functionals differ. Artificial daily labels enumerate sessions without providing exchange-calendar information.

We apply a positive numerical floor to keep risk forecasts scoreable and record its activations. Each comparator and corrected VolPFN seed also fits one QLIKE multiplier on 72 archived SPX, FTSE and N225 origins in 2012--2013. That multiplier is fixed before 2015--2016 evaluation and changes the level, rather than the ordering, of forecasts across dates. The calibration set covers only those three markets.

Calibration does not make all outputs equivalent: multiplying a median generally does not turn it into a conditional mean. The benchmark therefore evaluates calibrated pipelines rather than optimal interfaces. By contrast, the prior experiment uses raw moments, avoiding a fitted multiplier that could absorb a prior-related level difference. Raw-return quantiles support adapter checks, while the brief LSTM integration check contributes no substantive performance evidence.

\subsection{Common-period and realised-risk interfaces}

For the common-period study, intersecting fixed eligibility masks gives 22,026 market--date \emph{keys} in 2019--February 2022. A key identifies one market on one date. This alignment adds neither training nor calibration. Squared-return scoring uses second moment, while harmonised RK scoring uses central variance. Appendix~\ref{sec:common-period-units} records unit conversions and explains how this differs from the original RK diagnostic.

Information access also distinguishes the procedures. Standalone VolPFN uses daily returns; realised-risk helpers add lagged RK; and some competitors receive foreign histories. Combinations fit parameters on earlier completed observations pooled across markets. These information and estimation choices define the comparison. Ordering sessions alone verifies neither RK publication vintages nor equivalence between RK and session return risk.

\section{Uncertainty and decision rules}
\label{sec:inference}

We pair forecasts by market and date, average their scores across dates within each market, and then give markets and seeds equal weight. Longer histories therefore cannot dominate the panel result. Averaging individual model losses retains separately scored forecasts; it does not produce an ensemble prediction.

To retain persistent volatility and shared shocks when measuring date uncertainty, we resample calendar blocks with all markets moving together \citep{kunsch1989jackknife}. Missing sessions remain in place, and circular blocks wrap around the observed calendar. Block length approximates dependence under the stated procedure. It neither establishes stationarity nor provides distribution-free validity.

The prior study uses 5,000 resamples with sixty-date blocks and a one-sided test of improvement in the broad arm's common-$t_8$ score. Holm adjustment covers the frozen two-comparison family \citep{holm1979simple}, including the unrun broad-versus-GARCH slot at $p=1$. Support requires adjusted $p\leq0.05$, a positive basic 95\% interval, an observed mean of at least 0.005 nats and positive means in at least fifteen of twenty markets. These requirements address uncertainty, effect size and breadth separately. The size and breadth thresholds are protocol decisions rather than universal literature cutoffs.

Equivalence requires separate evidence that the difference is sufficiently small. The two one-sided tests require the basic 90\% interval to lie strictly within $\pm0.003$ nats. Failing to detect superiority does not satisfy this requirement. The superiority and equivalence endpoints have no joint familywise guarantee.

The general benchmark follows its own frozen protocol: 20,000 resamples, ten-date blocks, two-sided centred tests and Holm adjustment over three comparisons. Bonferroni basic intervals give approximate 95\% family coverage, while five- and twenty-date sensitivities are corrected separately. Neither the benchmark's ten-date length nor the prior study's sixty-date length is an estimated optimum or selected for significance. The benchmark adds no effect-size or breadth gate (Appendix~\ref{app:closure_protocol}).

All intervals condition on the fitted models, markets, scales and eligibility masks. They omit future training runs, calibration-estimation uncertainty and research selection. Although we report individual seeds, the resampling unit is the date block. A small $p$-value therefore cannot be interpreted as a probability of universal superiority.

\paragraph{Separate inference for the added studies.}
For the component and financial neural studies, we use 20,000 synchronised calendar-block resamples. Primary blocks contain sixty dates, with twenty- and 120-date alternatives checking sensitivity to that choice. The tests are two-sided and centred, with a plus-one correction for finite resampling. We apply Holm adjustment separately to the three- and four-comparison families, retaining unavailable RECH's slot. Full remains the pre-specified external reference even though noaux attains a lower mean.

Separate families of 28, 38 and twelve comparisons govern the harmonised daily, intraday and cross-market-helper studies. Their complete rule requires passing adjusted tests, positive directional family intervals, both risk targets and positive gains for every reference seed. We distinguish ordinary intervals from family intervals and report sensitivity families separately. Appendix~\ref{sec:new-research-results} retains every primary and unavailable contrast. Separate adjustments provide no dissertation-wide guarantee, and market signs describe consistency rather than independent replications.

\paragraph{Inference for the three completed follow-ups.}
The October follow-ups use endpoints suited to their separate questions. Historical and ETF comparisons test squared-return QLIKE, while native and common-$t_8$ density losses remain descriptive. The TTM controls test both risk targets. Each follow-up uses 20,000 synchronised circular calendar resamples, with sixty-date primary blocks and twenty- and 120-date sensitivities. Two-sided centred tests include the plus-one correction, and Holm adjustment is applied separately to two historical, four calibration and three ETF contrasts.

A basic marginal 95\% interval describes uncertainty for one contrast. The Bonferroni family intervals account for considering the declared set of contrasts together, so a marginal interval can exclude zero while its family interval includes it. We report both. Complete support requires a positive mean, adjusted $p<0.05$, a positive family interval and positive reference-seed gains; Appendix~\ref{sec:followup-results} retains the exact contrasts and sensitivities. Adjustment within these families does not provide a dissertation-wide guarantee. The intervals also condition on saved forecasts, omitting retraining, online-fit re-estimation and research-selection uncertainty.

\section{Exposure and scope}
\label{sec:exposure}

The corrected plans preceded corrected training and scoring, but earlier outcomes were already known. Later SPX, FTSE, GDAXI and N225 histories had prior audit contact, and an interim benchmark inspection covered 8,552 observations. We therefore report disclosed confirmations rather than globally untouched tests. Possible overlap with external models' pretraining data remains unknown.

The historical studies keep their original checkpoint and masks: 26,293 development observations in 2011--2014 and 50,073 later observations in 2015--February 2022. Tail procedures retain their numerical safeguards, while adaptive-tail evidence remains exploratory. Recording identities and eligibility masks makes these boundaries traceable without removing the earlier exposure.

The added studies also revisit exposed panels. Aligning dates changes their comparison mask, not their exposure, and the intraday and later helper extensions remain retrospective. Historical training and stopping rules were fixed before running, but its analysis specification was recorded after the first new seed result was visible and before the other four completed. TTM controls revisit the same exposed panel. These documented retrospective analyses are not formal preregistrations.

Before ETF acquisition, the protocol froze the roster, endpoints, comparisons and original checkpoint identities. Execution subsequently classified the study as qualified external population validation. A targeted project inventory found no ETF outcomes, although prior human exposure remains unverified. Shared dates and shocks with the index studies prevent a claim of globally untouched, independent confirmation. Fixed weights establish transfer under the recorded protocol without establishing independence from the research history. Appendix~\ref{sec:followup-methods} preserves this chronology. Restricting inputs to completed observations prevents prediction-time leakage, but cannot remove research-selection bias.

Chapter~\ref{ch:forecast} evaluates the issued forecasts under these scope and exposure limits. Chapter~\ref{ch:analysis} then tests explanations of their performance through matched interventions.
